\sscp{\tsc{Minahasan}}{C-09-02}
The following languages are all Austronesian,
belonging to the \tsc{Minahasa} subgroup,
also classified by \citet{Blust2019} as a member of \tsc{Phillipine}.
\\

\noindent{\wg\st{2pt}\begin{tabularx}{\linewidth}
{lllll>{\raggedright\arraybackslash}X}
\lsptoprule
%Language	&	ISO	&	Term	&	Gloss	&	Etymon	&	Source	\\	\midrule
%\endfirsthead
Language	&	ISO	&	Term	&	Gloss	&	Etymon	&	Source	\\	\midrule
%\endhead
Tondano			&	tdn	&	kuse;						&	cuscus;				&	**kusay			&	\cite[151]{Sneddon1978}	\\	\hhline{~}
						&			&	tembuŋ					&	small kind 		&	\it{misc.}	&	\cite[161,]{Stokhof1983b}	\\	\hhline{~}
						&			&									&	of cuscus			&							&	174												\\	
Tontemboan	&	tnt	&	‹kusé›;					&	bear cuscus;	&	**kusay			&	\cite[151]{Sneddon1978}	\\	\hhline{~}
						&			&	‹tĕmbuŋ/tĕbuŋ›	&	dwarf cuscus	&	\it{misc.}	&	\cite[205, 497, 502]{Schwarz1908}	\\
Tombulu			&	tom	&	kuse						&	cuscus				&	**kusay			&	\cite[151]{Sneddon1978}	\\
Tonsea			&	txs	&	kuse						&	cuscus				&	**kusay			&	\cite[151]{Sneddon1978}	\\
\lspbottomrule\end{tabularx}}

%\noindent{\wg\st{2pt}\begin{xltabular}{\linewidth}
%{lllll>{\raggedright\arraybackslash}X}
%\lsptoprule
%%Language	&	ISO	&	Term	&	Gloss	&	Etymon	&	Source	\\	\midrule
%%\endfirsthead
%Language	&	ISO	&	Term	&	Gloss	&	Etymon	&	Source	\\	\midrule
%\endhead
%
%\lspbottomrule
%\end{xltabular}}